% META: {"id": "1NSI-ARCHOS-EX-001", "chapitre": "1NSI-ARCHITECTURE-OS", "type_objet": "exercice", "capacites": ["P-ARCH-01A"], "parcours": 1, "competences": ["analyser"], "duree_min": 10, "mode_creation": "ex_nihilo", "sources_inspiration": [], "fichier_tex": "chapitres/1NSI-ARCHITECTURE-OS/exercices/1NSI-ARCHOS-EX-001.tex", "corrige_tex": "chapitres/1NSI-ARCHITECTURE-OS/corriges/1NSI-ARCHOS-CO-001.tex", "status": "manual_review"}
\begin{exercice}{1NSI-ARCHOS-EX-001}{1}{10}
\begin{enumerate}
  \item Dans l'architecture de von Neumann, quel constituant exécute les instructions, et quel constituant les stocke (avec les données) ?
  \item Un ordinateur a $8$ cœurs. Peut-il exécuter plusieurs instructions de programmes différents \textbf{en même temps} ? Justifier en citant le vocabulaire du cours (mono/multiprocesseur).
  \item Pourquoi le processeur ne peut-il pas, en lisant une suite de bits en mémoire, savoir \og à l'avance \fg{} s'il s'agit d'une instruction ou d'une donnée ?
\end{enumerate}
\end{exercice}

% BEGIN-VERIFY
% # Q3 : une meme suite de bits peut se lire comme une instruction ou comme une
% # donnee ; rien dans le contenu ne tranche. On le montre sur un octet.
% octet = 0b01100001
% assert octet == 97
% assert chr(octet) == "a"                    # lu comme un caractere
% assert octet == 0x61                        # lu comme un entier hexadecimal
% # Le meme motif interprete comme un code d'instruction fictif.
% jeu_instructions = {0b01100001: "CHARGE"}
% assert jeu_instructions[octet] == "CHARGE"
% # Trois lectures, un seul contenu : c'est la PHASE du cycle qui decide.
% assert len({chr(octet), str(octet), jeu_instructions[octet]}) == 3
%
% # Q2 : huit coeurs permettent huit executions simultanees ; un seul coeur
% # n'en permet qu'une a la fois.
% coeurs = 8
% assert coeurs > 1                            # multiprocesseur
% assert coeurs >= 8
% END-VERIFY
